\documentclass[fleqn,10pt]{wlscirep}
\usepackage[utf8]{inputenc}
\usepackage[T1]{fontenc}
%\documentclass{article}
%\usepackage[a4paper, total={6.5in, 10in}]{geometry}
%\usepackage[utf8]{inputenc} % allow utf-8 input
%\usepackage[style=numeric,sorting=none]{biblatex}
%\addbibresource{bibliography.bib}
\usepackage{lmodern}
\usepackage{amsmath}
\usepackage{amsthm}
\usepackage{amsfonts}
\usepackage{amssymb}
\usepackage{amsmath,amscd}
\usepackage{physics}
\usepackage{dsfont}
\usepackage{empheq}
\usepackage{graphicx}
\usepackage{subfig}
\usepackage{tabularx}
\usepackage{nicefrac}
\usepackage{xcolor}
\usepackage{float}
\usepackage{caption}
\usepackage{braket}
\usepackage{algorithm}
\usepackage{algpseudocode}
%\usepackage{subcaption}
\usepackage{booktabs}
\usepackage[toc]{appendix}
\usepackage{multirow}
\usepackage{subcaption}
\captionsetup{compatibility=false}
\usepackage{xr}



\title{Leveraging Analog Quantum Computing with Neutral Atoms for Solvent Configuration Prediction in Drug Discovery}

\author[1, +]{Mauro D'Arcangelo}
\author[2, +]{Daniele Loco}
\author[1]{Fresnel team}
\author[2,3,4]{Nicolaï Gouraud}
\author[2]{Stanislas Angebault}
\author[2]{Jules Sueiro}
\author[3]{Pierre Monmarché}
\author[2]{Jérôme Forêt}
\author[1]{Louis-Paul Henry}
\author[1,*]{Loïc Henriet}
\author[2,4,*]{Jean-Philip Piquemal}

\affil[1]{Pasqal, 2 avenue Augustin Fresnel, 91120 Palaiseau, France}
\affil[2]{Qubit Pharmaceuticals, Advanced Research Department, 24 rue du Faubourg Saint-Jacques, 75014 Paris, France}
\affil[3]{Sorbonne Université, Laboratoire Jacques-Louis Lions, UMR 7589 CNRS, 75005, Paris, France}
\affil[4]{Sorbonne Université, Laboratoire de Chimie Théorique, UMR 7616 CNRS, 75005, Paris, France}
\affil[*]{loic.henriet@pasqal.com, jean-philip.piquemal@sorbonne-universite.fr}

\affil[+]{these authors contributed equally to this work}

\keywords{Keyword1, Keyword2, Keyword3}

%----Helper code for dealing with external references----
% (by cyberSingularity at http://tex.stackexchange.com/a/69832/226)

\usepackage{xr}
\makeatletter

\newcommand*{\addFileDependency}[1]{% argument=file name and extension
\typeout{(#1)}% latexmk will find this if $recorder=0
% however, in that case, it will ignore #1 if it is a .aux or 
% .pdf file etc and it exists! If it doesn't exist, it will appear 
% in the list of dependents regardless)
%
% Write the following if you want it to appear in \listfiles 
% --- although not really necessary and latexmk doesn't use this
%
\@addtofilelist{#1}
%
% latexmk will find this message if #1 doesn't exist (yet)
\IfFileExists{#1}{}{\typeout{No file #1.}}
}\makeatother

\newcommand*{\myexternaldocument}[1]{%
\externaldocument{#1}%
\addFileDependency{#1.tex}%
\addFileDependency{#1.aux}%
}
%------------End of helper code--------------

% put all the external documents here!
\myexternaldocument{main}

\begin{abstract}
  We introduce quantum algorithms able to sample equilibrium water solvent molecules configurations within proteins thanks to analog quantum computing. To do so, we combine the 3D Reference Interaction Site Model (3D-RISM), an approach capable of predicting continuous solvent distributions to a quantum placement strategy. The intrinsic quantum nature of such coupling guarantees molecules not to be placed too close from each other’s, a constraint usually imposed by hand in classical approaches. We present first a full quantum adiabatic evolution model that uses a local Rydberg Hamiltonian to cast the general problem into an anti-ferromagnetic Ising model.  Its solution, an NP-hard problem in classical computing, is embodied into a Rydberg atoms array Quantum Processor Unit (QPU). Following a classical emulator implementation, a QPU portage allows to experimentally validate the algorithm performances on an actual quantum computer. As a perspective of use on next generation devices, we emulate a second hybrid quantum-classical version of the algorithm. Such a variational quantum approach (VQA) uses a classical Bayesian minimization routine to find the optimal laser parameters. Overall, these Quantum-3D-RISM (Q-3D-RISM) algorithms opens a new route towards the application of analog quantum computing in molecular modelling and drug design.  
\end{abstract}

\begin{document}

\flushbottom
\maketitle
\section*{Supplementary Information}
\section{Numerical simulation on CPUs for model systems with a QAE}
\label{sec:num-sim-supportQAE}

We briefly recall here that, given an array of $M$ quits, the result produced by either the VQA or the QAE algorithms is the position of $N$ molecules placed in correspondence to $N$ of the $M$ excited qubits. The results is coherent with the Rydberg blockade constraint and is the best at reproducing the reference density distribution. 

The qubit array is initially prepared on a 2D plane, as if one was to prepare it for an experimental set-up (see panel a, Figure~\ref{fig:exampleD}). Ideally, the origin of the reference system is the focus of the laser that one would use in a real experiment. This initial qubit array is then reduced as shown in Figure~\ref{fig:exampleD}, (panel b) based on a local density threshold. The 2D density is represented on the same plane where the qubits have beed placed initially. This procedure is used to reduce the number of qubits to emulate, so to dramatically reduce the computational cost.

\begin{figure}[!htbp]
\begin{center}
\subfloat[]{\includegraphics[width=0.40\textwidth]{img/2D/initial_qubit_registry.png}}
\subfloat[]{\includegraphics[width=0.59\textwidth]{img/2D/example_density.png}}
\end{center}
    \caption{(a) 2D qubits array used for the tests on the emulator; (b)  example of 2D density generated as a sum of gaussians; three gaussian distributions are summed up, and the respective centers are marked with red dots. Blue dots represents the only qubits, from the initial array in (a), which are placed in a region where the density is sufficiently high. The axes of the graph are spatial coordinates, in $\mu$m, while the colour palette from deep blue to yellow is used to mark the density level; level curves are also drawn on the same graph.}
    \label{fig:exampleD}
\end{figure}


The quantum adiabatic evolution (QAE) algorithm cannot be implemented on the current generation of neutral atom QPUs because of the lack of local addressing. As stated already in the main text, in other words, instead of the local terms $\Omega_i(t)$ and $\Delta_i(t)$, one only has access to a global control field $\Omega(t)$ and $\Delta(t)$. On the other hand, the local algorithm can be tested using a classical emulator. 

For these and the other emulations performed on a classical CPU (\textit{vide infra}) we used an AMD EPYC 7742 64-Core Processor CPU (3368.100 MHz).

To this aim we used the two 2D densities reported in Fig.~\ref{fig:QAE_2D_dens} (a and b subfigures).  We will refer in the following to those densities as \textit{synthetic densities}, since we generated them as a sum of gaussian distributions with known components. For each of case 16 qubits are used, placing them uniformely on the 2D densities.
The output of each emulation is essentially a sampling of the optimized $N$-qubits state $\ket{\Psi^{\{\Omega_{i}(t'),\Delta_{i}(t')\}}}$. This amount to sample the different qubits bitstrings, differing from total number of excitations and excited qubits, composing the state. The corresponding states obtained for the two test systems are reported in ~\ref{fig:QAE_2D_dens} c and d. The probability of those states best representing the gaussian component location is highlighted in red, in panel c and d.
\begin{figure}[!htbp]
\begin{center}
\subfloat[]{\includegraphics[width=0.5\textwidth]{img/Ising_2D/QA/2G.png}}
\subfloat[]{\includegraphics[width=0.5\textwidth]{img/Ising_2D/QA/3G.png}} \\
\subfloat[]{\includegraphics[width=0.5\textwidth]{img/Ising_2D/QA/2G_sampling.png}}
\subfloat[]{\includegraphics[width=0.5\textwidth]{img/Ising_2D/QA/3G_sampling.png}}
\end{center}
    \caption{Representation of the test synthetic sum of gaussian densities in 2D used to the the local laser based algorithm. The density maps, with level curves, are represented together with the qubits array (blue dots) built to address the location of the gaussian components. The centroid of each gaussian component is highlighted in red}
    \label{fig:QAE_2D_dens}
\end{figure}

From these tests we can observe how, between the $\sim 2^{16}$ possible states composing the $N$-qubits wavefunction, the QAE algorithm is able to find, with a good probability, the bitstring corresponding to the right placement of the gaussian distributions in each mixture.





\section{Numerical simulation on CPUs for model systems with a VQA}
\label{num-sim-support}

For the VQA algorithm, a bayesian optimizer is employed to optimise the laser pulses and solve the posed problem. Several copies of the emulation of a given system are performed. We test the four different minimizers\footnote{Baysian Minimization functions are provided by the the \texttt{skopt} python library and used in the Pulser software} namely: i) Dummy, ii) Forest, iii) GBRT and iv) GP minimizer.

The parameters space size, denoted $m$, of $\Omega(t)$ and $\Delta(t)$ is  arbitrary. The larger is the 2$m$ dimensional parameter space, the larger is the flexibility left to the shape of the laser pulse. We tested $m$ values between $m = 3$ and $m = 25$. 

For each emulated system, 200 bayesian optimization cycles are performed. At the end of each cycle $k$, the reconstruction of the quatum state $\ket{\Psi^{\Omega,\Delta}_k}$ is emulated, so to evaluate the quality of the newly proposed set of laser parameters $(\Omega_k(t), \Delta_k(t))$, according to the cost function used by one of the proposed algorithms. The sampling procedure is emulated through Pulser and ideally corresponds to performing $200$ measures on $\ket{\Psi^{\Omega,\Delta}_k}$ on the real QPU. This number of repeated measures is a realistic parameter, affordable also on the real QPU. This corresponds in practice to sampling the frequency of appearence of the basis states in the $N$-qubits wavefunction
\begin{equation}
    \ket{\Psi^{\Omega,\Delta}_k} = \sum_{i=1}^{2^N} a_i\ket{e_i}.
\end{equation}

Assuming one, or several best configurations of excited atoms can be defined, namely the best basis state $\ket{e^{*}_{l}}$ in eq~\ref{eq:bs_best} in the main text, we measure the performance of the algorithm from the value of $||a_l^{*} ||^2$. In practice this is measured as the frequency of the state $\ket{e^{*}_{l}}$ in the sample obtained with 200 shots on the $N$-qubits quantum state applying the optimized lase pulse $(\Omega(t)^*,\Delta(t)^*)$.

For emulation purposes, we need to keep the size of the system under 15-18 qubits so to reduce the cost of solving numerically the TDSE, and allow us to perform a larger number of optimization cycles. 

\subsection{Insights on the quantum evolution}
A set of $M$-qubits wavefunctions is generated during the evolution, each represented as a superposition of qubits basis bitstrings, according to their respective frequencies. To do so, the quantum evolution process is ran (or emulated), starting from a state $\ket{\psi_k}$, to end up in the final state $\ket{\psi_{T}}$, where $T$ is the evolution period of the quantum state.
In other words, starting from $k=0$ and a prepared initial state $\ket{\psi_0}$, the Schrödinger equation is solved to yield $\ket{\psi_{k+1}}=SolveSchrodinger(\ket{ \psi_{k}})\;\forall k=0,\dots,(T-1)$ (cf. Algorithm~\ref{alg:HQC}).

 
The quantum evolution is repeated $n_c=200$ times, in as many Bayesian optimization cycles, yielding $\ket{\psi_T}_i$ for all $i=1,\dots,200$. For every generated $\ket{\psi_T}_i$ every accessible basis states $\ket{e}_i$ are sampled and their occurrences stored in $freq_{\ket{e}}$ yielding the aforementioned expected result. The sampling procedure is emulated through Pulser and  corresponds to performing $n_m:=200$ measures on the real QPU.

\subsection{Insights on the Baysian optimizer}
Concerning the proposed VQA, a Bayesian optimizer is used to guide the choice of the laser parameters. For each cycle of the Bayesian optimization, the best form of the functions $\Omega(t)$ and $\Delta(t)$ is searched. The quality of the explored functions shapes is measured by the distance between the reference density and the density reproduced by the qubits quantum state $\ket{\psi}$.

The parameters space size for the laser pulse optimization, namely $\Omega(t)$ and $\Delta(t)$, is itself an arbitrary parameter required by the algorithm. The two functions are described by a sum of elementary functions: $\sum_{i=1}^mf_m(t)$. The upper bound $m\in\mathbb{N}$ of the sum corresponds to the size of the parameters space.

\subsection{Numerical tests}

To test the performances of the variational algorithm described in Algorithm~\ref{alg:HQC_IsG}, we used three simple 2D synthetic densities, reported in Figure~\ref{fig:2D_dens}. 

We have chosen on purpose simple synthetic densities to limit the number of qubits to employ, making possible to perform multiple times, during the optimization cycles, the numerically expensive emulations on classical CPUs. The Pulser emulator is used for this purpose, running the calculations on Atlas, the Qubit's hybrid CPU-GPU proprietary cluster. The best configuration of excited qubits is in these cases easy to identify: the set of excited qubits has to be composed by the qubits closer to the gaussians centroids, highlighted in red in Fig.~\ref{fig:exampleD}. A full cycle of optimization is performed for each test system, with the role of the QPU is taken by a classical solver for the time-dependent Shr{\"o}dinger equation.

For the second (VQA) proposed algorithm, the output of each emulation is the sampling of the optimized $N$-qubits state $\ket{\Psi^{\Omega^*,\Delta^*}}$, obtained after the hybrid quantum/classical optimization cycles. The optimizer employed employs Bayesian inference, introducing a stocastic behavior into the optimization procedure. Instead of performing a single optimization, reporting the final sampling of $\ket{\Psi^{\Omega^*,\Delta^*}}$, we perform several independent emulations for each test system. 

\begin{figure}[!htbp]
\begin{center}
\subfloat[]{\includegraphics[width=0.5\textwidth]{img/Ising_2D/TC3_dens.png}}
\subfloat[]{\includegraphics[width=0.5\textwidth]{img/Ising_2D/TC5_dens.png}} \\
\subfloat[]{\includegraphics[width=0.5\textwidth]{img/Ising_2D/TC6_dens.png}}
\end{center}
    \caption{Representation of the test synthetic sum of gaussian densities in 2D. The density maps, with level curves, are represented together with the qubits array (blue dots) built to address the location of the gaussian components. The centroid of each gaussian component is highlighted in red}
    \label{fig:2D_dens}
\end{figure}


\begin{figure}[!htbp]
\begin{center}
\subfloat[]{\includegraphics[width=0.4\textwidth]{img/Ising_2D/TC3.png}}
\subfloat[]{\includegraphics[width=0.4\textwidth]{img/Ising_2D/TC5.png}} \\
\subfloat[]{\includegraphics[width=0.4\textwidth]{img/Ising_2D/TC6.png}}
\subfloat[]{\includegraphics[width=0.4\textwidth]{img/Ising_2D/TC6_ave_moreM.png}} 
\end{center}
    \caption{Population of the best qubit-string (BS), representing the solution for the three test densities; the population is reported as a percentage on the overall sampling of the optimized wave function $\ket{\Psi^{\Omega^*,\Delta^*}}$. Results from simulations using different algorithm depths are compared. The parameter m represents the size of the optimization space (number of points used to discretize the laser pulse shape). Higher the allowed parameter optimization is, the better the solution can be found by a suitable algorithm}
    \label{fig:2D_globalIsing}
\end{figure}

To have a global vision of the results, we report in Fig.~\ref{fig:2D_globalIsing} the averaged probability of only the best qubit-string (BS) $\ket{e^*}$, which best represents the reference density. For the four and two components gaussian densities (a and b in Fig.~\ref{fig:2D_globalIsing}), the BS represents between 100\% and 60\% of the total sample of the optimized state. Such an high percentage is also due to the fact that the actual solution of the problem corresponds to an MIS on the qubit array, which is the best scenario for Pasqal's QPU. Nevertheless, also the third case (c in Fig.~\ref{fig:2D_dens}), with a denser qubit distribution, shows a satisfactory accuracy, even if the contribution of $\ket{e^*}$ drops quite significantly. In this last case we show the beneficial effects of enlarging the allowed parameter space for the optimization, improving the level of detail for the laser pulse shaping (m parameter reported in all histograms).

\section{Numerical simulation on 3D-RISM 2D slices from MPU-I test system}

Here we report the emulation results, using the VQA, obtained for all the six 2D slices obtained from the 3D-RISM density computed from the MUP-I system, with co-crystallized ligand (2-sec-butyl-4,5-dihydrothiazole.
In Fig.~\ref{fig:3DRISM_slices_res} the six slices, and corresponding output hinstograms, are reported and compared.

%\begin{figure}[!htbp]
%\begin{center}
%\subfloat[]{\includegraphics[width=0.5\textwidth]{img/3DRISM/MUP/slices/40519-outputs-selected_qr_plus_density.png}}
%\subfloat[]{\includegraphics[width=0.5\textwidth]{img/3DRISM/MUP/slices/40515-outputs-selected_qr_plus_density.png}}\\
%\subfloat[]{\includegraphics[width=0.5\textwidth]{img/3DRISM/MUP/slices/40516-outputs-selected_qr_plus_density.png}}
%\subfloat[]{\includegraphics[width=0.5\textwidth]{img/3DRISM/MUP/slices/40517-outputs-selected_qr_plus_density.png}}\\
%\subfloat[]{\includegraphics[width=0.5\textwidth]{img/3DRISM/MUP/slices/40518-outputs-selected_qr_plus_density.png}}
%\end{center}
    %\caption{Additional 3D-RISM density slices, computed for the major urinary protein (MUP-I) and used as a practical test case. The density maps, with level curves, are represented together with the qubits array (blue dots)}
%    \label{fig:3DRISMs_slice_em}
%\end{figure}

\begin{figure}[!htbp]
\begin{center}
\includegraphics[width=0.99\textwidth]{img/3DRISM/MUP/slices/3DRISM_slices.png}
\end{center}
    \caption{Representation of the test densities computed for the MUP-I protein, including the co-crystalized ligand, using the 3D-RISM method. We reported five 2D slices of the total 3D density, with level curves, on which the qubits array (blue dots) is built, to address the location of the signle gaussian components. The centroid of each gaussian component, found as an algorithm solution, is highlighted in red, for the highest probability bitstring, and orange, for the second highest probability bitstring.}
    \label{fig:3DRISM_slices_res}
\end{figure}


\subsection*{A short digression on 3D-RISM}

Statistical mechanical theories, based on distribution functions and integral equations, can provide a rigorous framework for describing liquid phases, incorporating the granularity, packing, and specific molecular interactions without needing an explicit atoms-based representation of large ensembles of solvent molecules.

An example is the reference interaction site model (RISM)~\cite{chandler_optimized_1972SI,chandler_derivation_1973SI,beglov_integral_1997SI,perkyns_sitesite_1992,perkyns_dielectrically_1992,roy_chapter_2021,chandler_equilibrium_1974SI} proposed first by Chandler and Andersen in 1972, and 3D-RISM (the version of RISM extended to 3D)~\cite{cortis_three-dimensional_1997,kovalenko_three-dimensional_2003SI,kovalenko_self-consistent_1999,luchko_three-dimensional_2010,roy_chapter_2021}. 

3D-RISM represents a widely spread approach to treat the problem of solvation, as testified by the highly cited original works and by the implementation of the model in one of the standard software for molecular dynamics simulations, namely the Amber suite of programs.

The 3D-RISM theory provides probability density of all the solvent interaction sites $\gamma$ around a solute at position $\mathbf{r}$ as the product 
\begin{equation*}
    \rho_{\gamma}(\mathbf{r}) = \rho_{\gamma}g_{\gamma}(\mathbf{r})
\end{equation*}

of the average number density $\rho_{\gamma}$ in the bulk solution, and the
normalized density distribution $g_{\gamma}(\mathbf{r})$

The density distribution $g_{\gamma}(\mathbf{r})$ is the fundamental output of the model, with a site $\gamma$, corresponding to one of the physically different atoms which define the solvent model. 
High values of $g_{\gamma}(r)$ corresponds to regions with high solvent concentration, compared to the pure solvent bulk. One can interpret those high density regions as
location of stable water molecules, due to strong water-solute interactions.

In practice, 3D-RISM is an integral, site-site theory which needs as input the molecular structure of solute which one needs to solvate, and the solute and solvent models parameters, to compute the inter- and intra-molecular interactions between them.

A grid representation of the solvent around the solute, is used. The $g(r)$ function is evaluated on the grid, and each site of the solvent model will have a corresponding distribution value.

The 3D-RISM fundamental integral equation is the the solute-solvent total correlation function (TCF)~\cite{kovalenko_three-dimensional_2003SI,kovalenko_self-consistent_1999,luchko_three-dimensional_2010}

\begin{equation}
    h(r) = \sum_\alpha \int  c_\alpha^{UV}(r - r')\chi_{\alpha}^{VV}(r') dr^\prime
\end{equation}


for the solvent site $\gamma$. The superscripts ``U'' and ``V'' denote the solute and solvent species; the sum runs over the $\alpha$ solvent sites. $c_\alpha^{UV}$ is the 3D direct correlation function (DCF) for solvent site $\alpha$, which is proportional to the interaction potential between the solute and solvent site $u_\alpha^{UV}(r)/k_{B}T$, with opposite sign. $u_\alpha$ can be seen as the sum of pairwise solute-solvent, site-site potentials given from all the solute interaction sites located at frozen positions $R_i$.

The site-site susceptibility of the solvent is given by

\begin{equation}
\chi_{\alpha\gamma}^{VV}(r) = \omega_{\alpha\gamma}^{VV}(r) + \rho_\alpha h_{\alpha\gamma}^{VV}(r) 
\end{equation}

with $\omega_{\alpha\gamma}^{VV}$ the intramolecular correlation function,  representing the internal geometry of the solvent molecules, while $h_{\alpha\gamma}^{VV}$ is the site-site radial total correlation function of the pure solvent calculated from RISM~\cite{chandler_optimized_1972SI,beglov_integral_1997SI,perkyns_sitesite_1992,perkyns_dielectrically_1992}. A reference model for the solvent is chosen, and the susceptibility is computed accordingly. $\rho_\alpha$ is the solvent site bulk density.  

\begin{equation}
    h(r) = g(r) - 1
\end{equation}

is the relation which connects the main 3D-RISM integral equation to the desired solvent structure distribution function $g(r)$. 
Here, $h_\alpha$ and $c_\alpha$ are unknown. Therefore, in order to calculate $g$, the 3D-RISM algorithms uses an additional closure equation given by a Taylor expansion of order n:

\begin{equation}
    g_\alpha(r) = \left\{ \begin{array}{ll} \exp(t_\alpha(r)) & \text{if }  t_\alpha < 0 \\ \sum_{i=0}^{n} \frac{(t_\alpha(r))^i}{i!} & \text{if }  t_\alpha \geq 0 \end{array} \right.
\end{equation}

Where $t_\alpha(r) = -\frac{u_\alpha}{k_B T} + h_\alpha(r)-c_\alpha(r)$. When  $n=1$, this equation is called the Kovalenko-Hirata (KH) closure. When $n=\infty$, it is the so-called hypernetted $\lambda$-chain equation (HNC).

A lot of attention has been put in applying the model to molecular liquids~\cite{kinoshita_application_1996,roy_application_2020}, to study the solvation of
small molecules~\cite{truchon_cavity_2014,johnson_small_2016} or ion solvation~\cite{kovalenko_potentials_2000}, and even recently to larger biomolecules.~\cite{imai_ligand_2009,luchko_three-dimensional_2010,perkyns_protein_2010,sindhikara_analysis_2013}

To our knowledge, only recently 3D-RISM has been applied to study the hydration states of protein binding pockets~\cite{masters_efficient_2018,fusani_optimal_2018SI,nguyen_molecular_2019}, where classical computing technologies, based either on genetic algorithms (GAsol~\cite{masters_efficient_2018,fusani_optimal_2018SI}) and/or on molecular dynamics.~\cite{masters_efficient_2018,nguyen_molecular_2019} 



\bibliography{SI_biblio.bib} 

\end{document}